% !TeX program = xelatex
% ======================================================================
% 全国大学生数学建模竞赛（CUMCM）2026 论文模板
%
% 依据本次对话中已确认的 2026 官方要求搭建：
% 1. A4；页边距上下左右至少 2.5 cm；
% 2. 电子论文第一页为“论文标题 + 摘要 + 关键词”，摘要原则上不超过 1 页；
% 3. 不需要英文摘要；
% 4. 不设置目录；
% 5. 正文不超过 30 页，附录页数不限；
% 6. 电子论文不包含承诺书、编号页；
% 7. 摘要、正文、附录不得出现学校、队员姓名、赛区等身份信息；
% 8. 参考文献前设置“AI工具使用声明”；
% 9. 附录中列出支撑材料文件，并给出完整、可运行的主要源程序；
% 10. PDF 与支撑材料压缩包均需控制在赛事要求的大小范围内。
%
% 推荐：
%   latexmk -xelatex main.tex
%
% 推荐项目结构：
%   main.tex
%   cumcmthesis.cls           % 若使用官方/常用 cumcmthesis 模板
%   figures/
%   code/
%   data/
%   support/
%       AI工具使用详情.pdf
%
% 注意：
% - 本模板“章节名称”采用国赛优秀论文常见写法，并非官方强制章节名称。
% - 比赛正式提交前，请再次核对全国组委会及所在赛区/学校最新通知。
% ======================================================================

% ----------------------------------------------------------------------
% 1. 文档类与字体（独立、可稳定嵌入中文字体）
% ----------------------------------------------------------------------
% fontset=none：不依赖操作系统自带的宋体/黑体等字体；
% 明确使用 Noto CJK 字体，避免 PDF 中“中文消失，只剩数字/斜杠”的问题。
\documentclass[UTF8,fontset=windows,a4paper,12pt]{ctexart}

\usepackage[
  a4paper,
  left=2.5cm,
  right=2.5cm,
  top=2.5cm,
  bottom=2.5cm
]{geometry}
\usepackage{amsmath,amssymb,bm}
\usepackage{graphicx}
\usepackage{booktabs}
\usepackage{longtable}
\usepackage{multirow}
\usepackage{listings}
\usepackage{xcolor}
\usepackage{subcaption}
\usepackage{float}
\usepackage[hidelinks]{hyperref}

% ---------- 字体体系：正文宋体风格 + 标题黑体风格 + 英文 Times 风格 ----------
% 中文：
%   正文  -> Noto Serif CJK SC（宋体/明朝体风格）
%   标题  -> Noto Sans CJK SC（黑体风格）
%   代码  -> Noto Sans Mono CJK SC
% 西文：
%   正文  -> Liberation Serif（Times New Roman 兼容风格）
%   标题  -> Liberation Sans
%   代码  -> Liberation Mono
%
% 所有字体均由 XeLaTeX 嵌入 PDF，避免浏览器预览时中文消失。
% 使用 ctex 的 Windows 字体方案：中文采用宋体/黑体/仿宋/楷体，
% 西文采用 Windows 可用字体，避免依赖未安装的 Liberation/Noto 字体。

\setlength{\parindent}{2em}
\linespread{1.18}

% 页码从摘要页开始，页脚居中
\pagestyle{plain}
\pagenumbering{arabic}
\setcounter{page}{1}

% “一、问题重述”“二、问题分析”……
\ctexset{
  section = {
    name = {,、},
    number = \chinese{section},
    format = \centering\sffamily\bfseries\zihao{4},
    aftername = \hspace{0.4em},
    beforeskip = 1.2ex plus .2ex,
    afterskip = 1.2ex plus .2ex
  },
  subsection = {
    format = \sffamily\bfseries\zihao{-4},
    beforeskip = 1.0ex plus .2ex,
    afterskip = .7ex plus .2ex
  },
  subsubsection = {
    format = \sffamily\bfseries\zihao{-4},
    beforeskip = .8ex plus .2ex,
    afterskip = .5ex plus .2ex
  }
}

\graphicspath{{figures/}}
\renewcommand{\arraystretch}{1.15}
\setlength{\tabcolsep}{6pt}

% ----------------------------------------------------------------------
% 3. 代码样式
% ----------------------------------------------------------------------
\lstset{
  basicstyle=\ttfamily\small,
  numbers=left,
  numberstyle=\scriptsize,
  numbersep=8pt,
  frame=single,
  breaklines=true,
  columns=fullflexible,
  showstringspaces=false,
  tabsize=4
}

% ----------------------------------------------------------------------
% 4. 论文标题（电子版匿名）
% ----------------------------------------------------------------------
\title{\sffamily\bfseries\zihao{2} XXXX问题的数学模型}
\author{}
\date{}
\providecommand{\keywords}[1]{\par\vspace{0.5em}\noindent\textbf{关键词：}#1}

% ----------------------------------------------------------------------
% 5. 常用数学命令
% ----------------------------------------------------------------------
\newcommand{\R}{\mathbb{R}}
\newcommand{\E}{\mathbb{E}}
\newcommand{\Var}{\operatorname{Var}}
\newcommand{\argmin}{\operatorname*{arg\,min}}
\newcommand{\argmax}{\operatorname*{arg\,max}}

\begin{document}

% ----------------------------------------------------------------------
% 版式说明（比赛写作时无需改）
% ----------------------------------------------------------------------
% 本模板采用：
% - 中文正文：宋体风格；
% - 中文标题：黑体风格；
% - 英文/数字：Times New Roman 兼容风格；
% - 数学公式：LaTeX 标准数学字体。
%
% 全国组委会对“字号、字体、行距、颜色”没有统一强制规定；
% 因此这里采用国内科技论文常见、阅读性较好的统一字体体系。
% ----------------------------------------------------------------------

% ======================================================================
% 第一页：论文标题 + 摘要 + 关键词
% 官方要求电子论文第一页直接进入摘要页，不含承诺书/编号页。
% 摘要原则上控制在 1 页以内；不写英文摘要。
% ======================================================================
\maketitle

\begin{abstract}
% ---------------- 摘要写作建议 ----------------
% 建议直接回答四件事：
% 1. 研究了什么；
% 2. 各小问分别建立了什么模型、用了什么方法；
% 3. 得到了哪些关键数值/结论；
% 4. 如何检验模型、模型价值是什么。
%
% 避免写成“本文首先……其次……最后……”的目录式流水账。
% 尽量给出关键指标、误差、最优参数、预测结果等定量信息。
% ----------------------------------------------

针对题目所研究的\textbf{核心对象与目标}，本文从\textbf{主要矛盾/关键机制}出发，
建立了由\textbf{模型一}、\textbf{模型二}和\textbf{综合优化模型}组成的统一分析框架。

针对问题一，首先对原始数据进行\textbf{清洗、重构与特征提取}，在若干合理假设下建立
\textbf{模型名称}。通过\textbf{求解算法}得到……，关键结果为……。

针对问题二，在问题一模型基础上引入\textbf{新增因素/约束}，建立\textbf{模型名称}，
并利用\textbf{算法或数值方法}求解。结果表明……。

针对问题三，进一步考虑\textbf{泛化场景/不确定性/优化目标}，构建\textbf{综合模型}。
通过\textbf{敏感性分析、误差分析或稳健性检验}验证模型可靠性，最终得到……。

综上，本文模型具有\textbf{可解释性、稳定性与一定推广能力}，可为……提供定量依据。

\keywords{关键词1\quad 关键词2\quad 关键词3\quad 关键词4}
\end{abstract}

% 官方要求：不设置目录
% \tableofcontents

% ======================================================================
\section{问题重述}
% ======================================================================

\subsection{问题背景}

“板凳龙”是一种由多节板凳首尾相连而成的传统舞龙道具。舞龙时，
各节板凳通过把手相互连接，龙头的运动会经由刚性连接逐节传递至龙身和龙尾。
题目所研究的板凳龙由 $1$ 节龙头、$221$ 节龙身和 $1$ 节龙尾组成，共计
$223$ 节板凳。龙头板长为 $3.41\,\mathrm{m}$，龙身和龙尾板长均为
$2.20\,\mathrm{m}$，所有板凳宽度均为 $0.30\,\mathrm{m}$；把手中心位于
板凳两端内侧 $0.275\,\mathrm{m}$ 处。相邻板凳在把手处铰接，因此整条
板凳龙可视为由定长杆件连接而成的平面刚性链。

为使板凳龙在有限场地内完成盘入、调头和盘出，需要合理确定盘入螺线、
调头曲线及其运动速度。题目要求在给定几何尺寸与运动条件下，计算各把手
的位置和速度，判断板凳之间是否发生碰撞，并进一步优化螺距、调头路径和
龙头行进速度。

% 写作提示：
% - 只保留与建模有关的背景；
% - 不整段照抄题面；
% - 通常 1--2 段即可。

\subsection{问题提出}

根据题意，本文需要解决以下问题：

\textbf{问题一：} 龙头前把手从盘入螺线第 $16$ 圈的 $A$ 点出发，以
$1\,\mathrm{m/s}$ 的速度沿螺距为 $0.55\,\mathrm{m}$ 的等距螺线顺时针盘入。
建立运动学模型，计算从 $0\,\mathrm{s}$ 到 $300\,\mathrm{s}$ 每秒整条板凳龙
各把手中心的位置和速度，并按题目要求给出典型时刻的计算结果及完整结果文件。

\textbf{问题二：} 在问题一盘入过程的基础上考虑板凳的实际宽度，建立板凳间的
碰撞判定模型，确定板凳龙继续盘入时首次发生碰撞前的终止时刻，并给出该时刻
所有把手中心的位置和速度。

\textbf{问题三：} 调头区域为以螺线中心为圆心、直径为 $9\,\mathrm{m}$ 的圆形区域。
在保证龙头能够沿螺线盘入至调头区域边界且途中不发生碰撞的条件下，确定盘入
螺线的最小螺距。

\textbf{问题四：} 当盘入螺线的螺距为 $1.7\,\mathrm{m}$ 时，以两段彼此相切且
分别与盘入、盘出螺线相切的圆弧构造 S 形调头曲线。在前一段圆弧半径为后一段
两倍的基础上，分析是否可以缩短调头路径，并计算调头前后 $-100\,\mathrm{s}$
至 $100\,\mathrm{s}$ 内各把手中心的位置和速度。

\textbf{问题五：} 在问题四确定的运动路径上，所有把手的速度均不得超过
$2\,\mathrm{m/s}$。确定满足该速度约束时龙头前把手所能达到的最大行进速度。

% 若题目只有三问，删除问题四及后续对应章节。

% ======================================================================
\section{问题分析}
% ======================================================================

\subsection{问题一分析}

问题一要求由龙头前把手的已知运动反推出其余 $223$ 个把手在各时刻的位置与
速度，其本质是等距螺线上的刚性链运动学问题。板凳龙共有 $224$ 个把手中心点，
相邻把手间距由相应板凳的长度和孔位共同确定；在盘入阶段，各把手中心均位于
同一条等距螺线上。因此，整条板凳龙的运动可以由龙头前把手的运动状态和相邻
把手间的定长约束逐节确定。

首先建立以螺线中心为原点的平面直角坐标系，将等距螺线参数化。由于龙头前把手
沿曲线以恒定速率运动，可由螺线弧长与时间的关系反求其极角和坐标。随后，以前一
把手为圆心、相邻把手的固定距离为半径，求该圆与螺线在外侧方向上的第一个交点，
从而依次确定其余把手的位置。最后，对相邻把手间的定长约束关于时间求导，建立
角速度递推关系，并据此计算各把手的瞬时线速度。对 $t=0,1,\ldots,300\,\mathrm{s}$
逐时求解，即可得到题目要求的完整位置和速度序列。

\subsection{问题二分析}

问题二在问题一基础上增加了……约束/变量，使得原问题由……
转化为……问题。为此，需要在前述模型中引入……，
并采用……方法对参数或决策变量进行求解。

\subsection{问题三分析}

问题三需要同时兼顾……、……与……，具有明显的多目标/非线性/不确定性特征。
因此本文拟构建……指标体系，并采用……方法统一衡量各因素，
最后通过……得到综合最优结果。

\subsection{问题四分析}

问题四重点考察模型的泛化能力。本文将在原模型的基础上调整……参数，
分析模型在不同……条件下的变化，并从……角度提出可操作建议。

\subsection{建模思路}

本文总体建模思路如下：
\begin{enumerate}
  \item 对题目数据进行清洗、单位统一、异常值检测与必要的特征提取；
  \item 针对问题一建立基础模型，确定关键变量与参数；
  \item 在基础模型上逐步加入问题二、问题三中的新因素和约束；
  \item 通过数值算法完成模型求解，并利用图表展示结果；
  \item 采用误差分析、敏感性分析或稳健性检验评价模型；
  \item 根据结果给出结论，并讨论模型的推广价值。
\end{enumerate}

% 将下方占位框替换为流程图：
\begin{figure}[H]
  \centering
  \fbox{\rule{0pt}{5.0cm}\rule{0.86\textwidth}{0pt}}
  \caption{总体建模流程图}
  \label{fig:workflow}
\end{figure}

% ======================================================================
\section{模型假设}
% ======================================================================

为了突出问题的主要矛盾并保证模型可求解，作如下假设：

\begin{enumerate}
  \item \textbf{假设一：} 题目所给数据真实、可靠，能够反映研究对象的基本规律；
  \item \textbf{假设二：} 在研究时间或空间范围内，未明确给出的外部因素对结果影响较小；
  \item \textbf{假设三：} 模型涉及的主要参数在相应研究区间内保持稳定；
  \item \textbf{假设四：} 对个别缺失或异常数据进行合理处理后，不改变整体变化趋势。
\end{enumerate}

% 正式比赛时必须改成与具体题目强相关的假设。
% 不建议为了凑数量写“空气阻力忽略”等与实际模型无关的套话。

% ======================================================================
\clearpage
\section{符号说明}
% ======================================================================

本文主要符号及其含义如表~\ref{tab:symbols} 所示。

\begin{table}[H]
  \centering
  \caption{主要符号说明}
  \label{tab:symbols}
  \begin{tabular}{cp{0.58\textwidth}c}
    \toprule
    符号 & 含义 & 单位 \\
    \midrule
    $t$ & 龙头前把手开始盘入后经过的时间 & $\mathrm{s}$ \\
    $p$ & 等距螺线的螺距 & $\mathrm{m}$ \\
    $b$ & 螺线参数，$b=p/(2\pi)$ & $\mathrm{m}$ \\
    $N$ & 板凳龙把手中心点的总数，$N=224$ & -- \\
    $i$ & 把手编号；$i=0$ 表示龙头前把手，$i=223$ 表示龙尾后把手 & -- \\
    $\theta_A$ & 龙头前把手在初始点 $A$ 处的极角，$\theta_A=32\pi$ & $\mathrm{rad}$ \\
    $\theta_i(t)$ & 时刻 $t$ 第 $i$ 个把手中心在螺线上的极角 & $\mathrm{rad}$ \\
    $P(\theta)$ & 极角为 $\theta$ 的螺线点的坐标向量 & $\mathrm{m}$ \\
    $P_i(t)$ & 时刻 $t$ 第 $i$ 个把手中心的坐标向量 & $\mathrm{m}$ \\
    $d_i$ & 第 $i$ 与第 $i+1$ 个把手中心之间的固定距离 & $\mathrm{m}$ \\
    $S(\theta)$ & 从螺线参数原点至极角 $\theta$ 处的弧长原函数 & $\mathrm{m}$ \\
    $\dot\theta_i(t)$ & 时刻 $t$ 第 $i$ 个把手中心的角速度 & $\mathrm{rad/s}$ \\
    $v_i(t)$ & 时刻 $t$ 第 $i$ 个把手中心的线速度大小 & $\mathrm{m/s}$ \\
    \bottomrule
  \end{tabular}
\end{table}

% 建议只列高频、重要、容易混淆的符号。
% 简单且只出现一次的变量，可直接在公式后解释。

% ======================================================================
\section{数据预处理}
% ======================================================================

\subsection{数据清洗}

首先对题目所给数据进行完整性与合理性检查，主要包括：
\begin{enumerate}
  \item 检查缺失值、重复值及明显异常值；
  \item 统一不同数据表中的单位、时间尺度与空间尺度；
  \item 根据数据特征采用插值、平滑或异常值修正方法；
  \item 对后续建模所需变量进行重构与特征提取。
\end{enumerate}

\subsection{描述性统计}

对关键变量计算均值、标准差、极值及相关性，并通过可视化识别主要变化规律。

\begin{figure}[H]
  \centering
  \fbox{\rule{0pt}{4.5cm}\rule{0.78\textwidth}{0pt}}
  \caption{关键变量描述性分析示意图}
  \label{fig:eda}
\end{figure}

% 若题目本身不存在数据预处理，可删掉整个第五章，
% 后续章节编号会自动调整。

% ======================================================================
\section{模型的建立与求解}
% ======================================================================

% ----------------------------------------------------------------------
\subsection{问题一模型的建立与求解}

\subsubsection{盘入螺线及把手编号}

以盘入螺线中心 $O$ 为原点建立平面直角坐标系，以从 $O$ 指向初始点
$A$ 的方向为 $x$ 轴正方向，并规定逆时针方向为极角正方向。题目给定螺距
$p=0.55\,\mathrm{m}$，故等距螺线的极坐标方程为
\begin{equation}
  r=b\theta,\qquad b=\frac{p}{2\pi}=\frac{0.55}{2\pi}.
  \label{eq:q1-spiral}
\end{equation}
相应的直角坐标参数方程可写为
\begin{equation}
  P(\theta)
  =b\theta
  \begin{bmatrix}
    \cos\theta\\
    \sin\theta
  \end{bmatrix}.
  \label{eq:q1-coordinate}
\end{equation}

龙头前把手初始位于第 $16$ 圈的 $A$ 点，因而
\begin{equation}
  \theta_A=32\pi,\qquad
  P_0(0)=P(32\pi)=(8.8,0)^\mathsf{T}\ \mathrm{m}.
  \label{eq:q1-initial}
\end{equation}
将龙头前把手编号为 $0$，沿龙身方向依次编号，龙尾后把手编号为 $223$。
由于把手中心距离板凳端部均为 $0.275\,\mathrm{m}$，龙头两把手间距和其余
板凳两把手间距分别为
\begin{equation}
  d_i=
  \begin{cases}
    3.41-2\times0.275=2.86, & i=0,\\
    2.20-2\times0.275=1.65, & i=1,2,\ldots,222.
  \end{cases}
  \label{eq:q1-segment-length}
\end{equation}

\subsubsection{龙头前把手的位置模型}

由式~\eqref{eq:q1-coordinate} 可得螺线弧长微元
\begin{equation}
  \mathrm{d}s
  =\sqrt{r^2+\left(\frac{\mathrm{d}r}{\mathrm{d}\theta}\right)^2}
   \,\mathrm{d}\theta
  =b\sqrt{1+\theta^2}\,\mathrm{d}\theta.
  \label{eq:q1-arc-element}
\end{equation}
取其弧长原函数为
\begin{equation}
  S(\theta)
  =\frac{b}{2}
  \left[
    \theta\sqrt{1+\theta^2}
    +\operatorname{arsinh}\theta
  \right].
  \label{eq:q1-arc-length}
\end{equation}
龙头前把手以 $v_0=1\,\mathrm{m/s}$ 沿螺线顺时针向内运动，故其极角随时间
减小。时刻 $t$ 的极角 $\theta_0(t)$ 满足
\begin{equation}
  S(\theta_A)-S\bigl(\theta_0(t)\bigr)=v_0t=t.
  \label{eq:q1-head-theta}
\end{equation}
函数 $S(\theta)$ 在 $\theta\geq 0$ 上严格单调递增，因此式
~\eqref{eq:q1-head-theta} 在区间 $[0,\theta_A]$ 上有唯一解。利用二分法求得
$\theta_0(t)$ 后，将其代入式~\eqref{eq:q1-coordinate}，即可得到龙头前把手坐标
\begin{equation}
  P_0(t)=P\bigl(\theta_0(t)\bigr).
  \label{eq:q1-head-position}
\end{equation}

\subsubsection{各把手位置的递推模型}

盘入过程中所有把手中心均位于同一条螺线上，且板凳可视为刚体，因此任意相邻
两把手之间满足定长约束
\begin{equation}
  \left\|P_{i+1}(t)-P_i(t)\right\|=d_i,
  \qquad i=0,1,\ldots,222.
  \label{eq:q1-distance-constraint}
\end{equation}
设已知第 $i$ 个把手的极角 $\theta_i$，将
$P_i=P(\theta_i)$、$P_{i+1}=P(\theta_{i+1})$ 代入式
~\eqref{eq:q1-distance-constraint}，可得
\begin{equation}
  b^2\left[
    \theta_i^2+\theta_{i+1}^2
    -2\theta_i\theta_{i+1}
      \cos(\theta_{i+1}-\theta_i)
  \right]=d_i^2.
  \label{eq:q1-position-recursion}
\end{equation}

由于后续把手位于前一把手的螺线外侧，应满足
$\theta_{i+1}>\theta_i$。式~\eqref{eq:q1-position-recursion} 可能存在多个数值根，
本文从 $\theta_i$ 开始向外搜索，选取满足定长条件的第一个根，以保证所得点为
物理上相邻的把手位置。依次令 $i=0,1,\ldots,222$，即可递推得到时刻 $t$
全部 $224$ 个把手的极角及坐标。

\subsubsection{各把手速度的递推模型}

直接对离散时刻坐标作差会引入截断误差，因此本文对定长约束进行解析微分。
令
\begin{equation}
  \Delta P_i=P_{i+1}-P_i,
\end{equation}
则式~\eqref{eq:q1-distance-constraint} 等价于
$\Delta P_i^\mathsf{T}\Delta P_i=d_i^2$。对时间求导可得
\begin{equation}
  \Delta P_i^\mathsf{T}
  \left(\dot P_{i+1}-\dot P_i\right)=0.
  \label{eq:q1-differentiated-constraint}
\end{equation}
由参数方程
\begin{equation}
  P'(\theta)
  =b
  \begin{bmatrix}
    \cos\theta-\theta\sin\theta\\
    \sin\theta+\theta\cos\theta
  \end{bmatrix},
  \qquad
  \dot P_i=P'(\theta_i)\dot\theta_i,
  \label{eq:q1-derivative}
\end{equation}
将其代入式~\eqref{eq:q1-differentiated-constraint}，得到角速度递推关系
\begin{equation}
  \dot\theta_{i+1}
  =\frac{
    \Delta P_i^\mathsf{T}P'(\theta_i)
  }{
    \Delta P_i^\mathsf{T}P'(\theta_{i+1})
  }\dot\theta_i,
  \qquad i=0,1,\ldots,222.
  \label{eq:q1-angular-velocity}
\end{equation}
龙头前把手的初始递推量由其恒定线速度确定：
\begin{equation}
  \dot\theta_0(t)
  =-\frac{v_0}{b\sqrt{1+\theta_0^2(t)}}.
  \label{eq:q1-head-angular-velocity}
\end{equation}
负号表示极角随盘入过程减小。根据
$\|P'(\theta)\|=b\sqrt{1+\theta^2}$，第 $i$ 个把手的线速度大小为
\begin{equation}
  v_i(t)
  =b\sqrt{1+\theta_i^2(t)}\,
   \left|\dot\theta_i(t)\right|.
  \label{eq:q1-linear-velocity}
\end{equation}

\subsubsection{数值求解方法}

根据上述模型，问题一的数值求解步骤如下：
\begin{enumerate}
  \item 输入螺距、板凳尺寸、初始极角和龙头速度，并计算各相邻把手间距；
  \item 对每个整数时刻 $t=0,1,\ldots,300\,\mathrm{s}$，在
  $[0,\theta_A]$ 上利用二分法求解式~\eqref{eq:q1-head-theta}；
  \item 由式~\eqref{eq:q1-position-recursion} 从龙头向龙尾逐节搜索，得到全部
  把手的极角和坐标；
  \item 由式~\eqref{eq:q1-angular-velocity} 递推各把手角速度，再由式
  ~\eqref{eq:q1-linear-velocity} 计算线速度；
  \item 将全部时刻的位置和速度保留 $6$ 位小数，写入题目要求的
  \texttt{result1.xlsx} 文件。
\end{enumerate}
二分法的角度误差限取 $10^{-12}$，最大迭代次数取 $200$。由于弧长函数单调，
且每次均选取前一把手外侧的第一个定长交点，递推过程具有明确且唯一的物理含义。

\subsubsection{计算结果与分析}

表~\ref{tab:q1-position} 给出了题目指定把手在典型时刻的位置，表
~\ref{tab:q1-speed} 给出了相应速度。完整的 $0$--$300\,\mathrm{s}$ 逐秒结果见
支撑材料 \texttt{result1.xlsx}。

\begin{table}[H]
  \centering
  \caption{问题一典型时刻各把手中心的位置（单位：$\mathrm{m}$）}
  \label{tab:q1-position}
  \resizebox{\textwidth}{!}{%
  \begin{tabular}{crrrrrr}
    \toprule
    把手位置 & $0\,\mathrm{s}$ & $60\,\mathrm{s}$ & $120\,\mathrm{s}$ &
    $180\,\mathrm{s}$ & $240\,\mathrm{s}$ & $300\,\mathrm{s}$ \\
    \midrule
    龙头前把手 $x$ & 8.800000 & 5.799209 & -4.084887 & -2.963609 & 2.594494 & 4.420274 \\
    龙头前把手 $y$ & 0.000000 & -5.771092 & -6.304479 & 6.094780 & -5.356743 & 2.320429 \\
    第1节龙身前把手 $x$ & 8.363824 & 7.456758 & -1.445473 & -5.237118 & 4.821221 & 2.459489 \\
    第1节龙身前把手 $y$ & 2.826544 & -3.440399 & -7.405883 & 4.359627 & -3.561949 & 4.402476 \\
    第51节龙身前把手 $x$ & -9.518732 & -8.686317 & -5.543150 & 2.890455 & 5.980011 & -6.301346 \\
    第51节龙身前把手 $y$ & 1.341137 & 2.540108 & 6.377946 & 7.249289 & -3.827758 & 0.465829 \\
    第101节龙身前把手 $x$ & 2.913983 & 5.687116 & 5.361939 & 1.898794 & -4.917371 & -6.237722 \\
    第101节龙身前把手 $y$ & -9.918311 & -8.001384 & -7.557638 & -8.471614 & -6.379874 & 3.936008 \\
    第151节龙身前把手 $x$ & 10.861726 & 6.682311 & 2.388757 & 1.005154 & 2.965378 & 7.040740 \\
    第151节龙身前把手 $y$ & 1.828753 & 8.134544 & 9.727411 & 9.424751 & 8.399721 & 4.393013 \\
    第201节龙身前把手 $x$ & 4.555102 & -6.619664 & -10.627211 & -9.287720 & -7.457151 & -7.458662 \\
    第201节龙身前把手 $y$ & 10.725118 & 9.025570 & 1.359847 & -4.246673 & -6.180726 & -5.263384 \\
    龙尾后把手 $x$ & -5.305444 & 7.364557 & 10.974348 & 7.383896 & 3.241051 & 1.785033 \\
    龙尾后把手 $y$ & -10.676584 & -8.797992 & 0.843473 & 7.492371 & 9.469336 & 9.301164 \\
    \bottomrule
  \end{tabular}%
  }
\end{table}

\begin{table}[H]
  \centering
  \caption{问题一典型时刻各把手中心的速度（单位：$\mathrm{m/s}$）}
  \label{tab:q1-speed}
  \resizebox{\textwidth}{!}{%
  \begin{tabular}{crrrrrr}
    \toprule
    把手位置 & $0\,\mathrm{s}$ & $60\,\mathrm{s}$ & $120\,\mathrm{s}$ &
    $180\,\mathrm{s}$ & $240\,\mathrm{s}$ & $300\,\mathrm{s}$ \\
    \midrule
    龙头前把手 & 1.000000 & 1.000000 & 1.000000 & 1.000000 & 1.000000 & 1.000000 \\
    第1节龙身前把手 & 0.999971 & 0.999961 & 0.999945 & 0.999917 & 0.999859 & 0.999709 \\
    第51节龙身前把手 & 0.999742 & 0.999662 & 0.999538 & 0.999331 & 0.998941 & 0.998065 \\
    第101节龙身前把手 & 0.999575 & 0.999453 & 0.999269 & 0.998971 & 0.998435 & 0.997302 \\
    第151节龙身前把手 & 0.999448 & 0.999299 & 0.999078 & 0.998727 & 0.998115 & 0.996861 \\
    第201节龙身前把手 & 0.999348 & 0.999180 & 0.998935 & 0.998551 & 0.997894 & 0.996574 \\
    龙尾后把手 & 0.999311 & 0.999136 & 0.998883 & 0.998489 & 0.997816 & 0.996478 \\
    \bottomrule
  \end{tabular}%
  }
\end{table}

由表~\ref{tab:q1-position} 可见，随龙头持续盘入，其到螺线中心的距离逐渐减小，
而位于后部的龙身和龙尾仍分布在较外侧的螺线上，体现了运动沿刚性链逐节传递的
特征。表~\ref{tab:q1-speed} 表明，龙头速度始终保持为 $1\,\mathrm{m/s}$；其余
把手速度均略低于龙头速度，并总体上从龙头向龙尾缓慢减小。在
$300\,\mathrm{s}$ 时，龙尾后把手坐标为
$(1.785033,9.301164)\,\mathrm{m}$，速度为 $0.996478\,\mathrm{m/s}$。
上述结果与等距螺线曲率由外向内逐渐增大的几何特征相符。

% ----------------------------------------------------------------------
\subsection{问题二模型的建立与求解}

\subsubsection{模型建立}

问题二在问题一基础上进一步考虑……。
因此，引入变量或参数……，将原模型扩展为
\begin{equation}
  \bm{x}^{*}
  = \argmin_{\bm{x}\in\Omega}F(\bm{x};\bm{\theta}),
  \label{eq:q2-model}
\end{equation}
其中，$\Omega$ 为可行域，$\bm{\theta}$ 为模型参数。

\subsubsection{模型求解}

采用……方法求解，并说明搜索区间、步长、初值、迭代次数、终止条件等关键设置。

\subsubsection{结果分析}

给出图表与关键数值，并重点回答：
\textbf{“结果是多少、为什么会这样、是否符合实际规律”。}

% ----------------------------------------------------------------------
\subsection{问题三模型的建立与求解}

\subsubsection{模型建立}

根据问题三要求，在已有模型上进一步引入……因素。
若问题涉及多目标，可构造综合目标函数
\begin{equation}
  S_i=\sum_{j=1}^{K}w_jz_{ij},
  \qquad
  \sum_{j=1}^{K}w_j=1,\quad w_j\ge 0.
  \label{eq:q3-score}
\end{equation}

若不同指标量纲不同，应在加权前进行合理的无量纲化处理。

\subsubsection{模型求解与结果}

根据综合得分或优化结果得到……，并给出排序、最优参数或最终策略。

% ----------------------------------------------------------------------
\subsection{问题四模型的建立与求解}

% 若题目只有三问，删除本节。

\subsubsection{模型推广}

将前述模型推广到……条件下。
通过修改……参数、……边界条件和……约束，即可得到新场景下的模型。

\subsubsection{推广结果与建议}

结合推广结果，从……、……和……三个方面提出建议。

% ======================================================================
\section{模型检验与敏感性分析}
% ======================================================================

\subsection{误差分析}

若模型包含预测或拟合部分，可采用均方根误差
\begin{equation}
  \mathrm{RMSE}
  = \sqrt{
      \frac{1}{n}
      \sum_{i=1}^{n}
      (y_i-\hat y_i)^2
    },
  \label{eq:rmse}
\end{equation}
并结合 MAE、MAPE、$R^2$ 或残差图评价模型精度。

% 请根据题型选择合适指标，不能机械套用 RMSE。

\subsection{敏感性分析}

选择关键参数 $\theta$，在基准值 $\theta_0$ 附近进行扰动：
\begin{equation}
  \theta=(1+\delta)\theta_0,
  \qquad
  \delta\in[-\varepsilon,\varepsilon].
\end{equation}

观察核心输出与最优方案是否发生显著变化，从而判断模型对参数变化的敏感程度。

\subsection{稳健性检验}

可根据题型采用：
随机噪声扰动、Bootstrap 重采样、Monte Carlo 模拟、
不同初值重复求解或不同数据子集交叉验证等方法。

% ======================================================================
\section{模型评价与改进}
% ======================================================================

\subsection{模型优点}

\begin{enumerate}
  \item \textbf{针对性强：} 模型围绕题目核心变量建立，能够直接回答各小问；
  \item \textbf{层次清晰：} 后续模型在前一问基础上逐步扩展，逻辑连贯；
  \item \textbf{可解释性较好：} 主要参数具有明确的实际意义；
  \item \textbf{稳健性较好：} 在合理扰动下主要结论保持稳定。
\end{enumerate}

\subsection{模型不足}

\begin{enumerate}
  \item 为保证模型可求解，对部分现实因素进行了简化；
  \item 部分参数依赖有限数据估计，可能存在一定误差；
  \item 在更复杂场景或更大数据规模下，求解成本可能增加。
\end{enumerate}

\subsection{模型改进}

可从数据精度、模型机制、参数估计、优化算法和场景推广等方面进一步改进。

% ======================================================================
\section{结论}
% ======================================================================

本文围绕题目提出的若干问题建立了相应数学模型，主要结论如下：

\begin{enumerate}
  \item \textbf{问题一：} ……，最终得到……；
  \item \textbf{问题二：} ……，最终得到……；
  \item \textbf{问题三：} ……，最终得到……；
  \item \textbf{问题四：} ……，最终得到……。
\end{enumerate}

综合来看，所建立模型能够较好描述……规律，并可为……提供定量参考。

% ======================================================================
% AI 工具使用声明
% 官方要求：放在“参考文献之前”。
% 二选一保留；若实际使用 AI，请保留第一种，并在支撑材料中提交
% “AI工具使用详情.pdf”。
% ======================================================================

\section*{AI工具使用声明}
\addcontentsline{toc}{section}{AI工具使用声明}

% ---------- 使用过 AI 时保留以下版本 ----------
本参赛队在竞赛过程中使用了AI工具，主要用于
\textbf{语言表达辅助、程序调试、资料整理或其他实际使用环节}，
所有涉及建模、计算、结果与结论的内容均由参赛队成员进行人工审核与核验，
详细使用情况见支撑材料中的《AI工具使用详情.pdf》。

% ---------- 完全未使用 AI 时，删除上段并取消下一行注释 ----------
% 本参赛队在竞赛过程中未使用任何AI工具。

% ======================================================================
% 参考文献
% ======================================================================

\begin{thebibliography}{99}

\bibitem{ref1}
作者.
\newblock 文献题目[J].
\newblock 期刊名, 年份, 卷(期): 页码.

\bibitem{ref2}
作者.
\newblock 书名[M].
\newblock 出版地: 出版社, 年份.

\bibitem{ref3}
机构名称.
\newblock 数据集或网页名称[EB/OL].
\newblock 访问日期：2026-09-XX.

\end{thebibliography}

% ======================================================================
% 附录
% 官方要求：附录页数不限。
% 应包含支撑材料文件列表，以及建模过程中使用的完整、可运行源程序等。
% ======================================================================

\clearpage
\appendix

\section{支撑材料文件列表}

% 请与最终提交的 ZIP/RAR 实际内容保持完全一致。
% 官方题目直接提供的原始附件通常不必重复放入支撑材料。
\begin{longtable}{p{0.34\textwidth}p{0.58\textwidth}}
  \caption{支撑材料文件列表}\label{tab:support-files}\\
  \toprule
  文件名 & 内容说明 \\
  \midrule
  \endfirsthead

  \toprule
  文件名 & 内容说明 \\
  \midrule
  \endhead

  code/q1.py & 问题一完整求解程序 \\
  code/q2.py & 问题二完整求解程序 \\
  code/q3.py & 问题三完整求解程序 \\
  data/processed.csv & 建模过程中生成的处理后数据 \\
  figures/ & 论文中使用的主要结果图 \\
  AI工具使用详情.pdf & AI 工具使用情况、典型交互、采纳与人工核验说明 \\
  \bottomrule
\end{longtable}

% 若未使用 AI，应删除“AI工具使用详情.pdf”这一行。

\section{完整源程序}

% 官方要求的是“完整、可运行”程序。
% 最终比赛版本不要只贴伪代码或删减后的核心片段。
%
% 程序较长时，可按问题拆成多个附录：
%   附录B 问题一程序
%   附录C 问题二程序
%   ……
%
% 下面仅作为格式示例。

\begin{lstlisting}[language=Python,caption={问题一完整程序示例},label={code:q1}]
import numpy as np

def objective(x):
    return np.sum(np.asarray(x) ** 2)

def main():
    x = np.zeros(5)
    result = objective(x)
    print(result)

if __name__ == "__main__":
    main()
\end{lstlisting}

\section{补充计算结果}

% 可放：
% - 大尺寸参数表
% - 大量中间结果
% - 敏感性分析补充图
% - 详细推导
% - 正文因篇幅未展示但用于支撑结论的数据

\begin{table}[H]
  \centering
  \caption{补充参数表}
  \begin{tabular}{ccc}
    \toprule
    参数 & 取值 & 说明 \\
    \midrule
    $\theta_1$ & -- & -- \\
    $\theta_2$ & -- & -- \\
    \bottomrule
  \end{tabular}
\end{table}

% ======================================================================
% 提交前最终检查（仅注释，不会出现在论文中）
%
% [ ] 第一页是否为“标题 + 摘要 + 关键词”
% [ ] 摘要是否原则上控制在 1 页以内
% [ ] 是否删除英文摘要
% [ ] 是否没有目录
% [ ] 正文是否不超过 30 页
% [ ] 页边距是否均不少于 2.5 cm
% [ ] 是否从摘要页开始显示页码
% [ ] 是否完全删除学校、姓名、赛区等身份信息
% [ ] 是否没有承诺书和编号页
% [ ] 所有外部资料是否规范引用
% [ ] AI工具使用声明是否位于参考文献之前
% [ ] 若使用 AI，支撑材料中是否有“AI工具使用详情.pdf”
% [ ] 附录是否列出支撑材料文件
% [ ] 源程序是否完整、可运行，且输出与论文一致
% [ ] 支撑材料文件列表与 ZIP/RAR 实际内容是否一致
% [ ] PDF 文件大小是否符合最新要求
% [ ] 支撑材料压缩包大小是否符合最新要求
% [ ] 是否再次核对全国组委会及所在赛区最新通知
% ======================================================================

\end{document}
